\section{Introduction}

Build systems avoid work because they are told what the work depends on. \texttt{make} compares
timestamps against a declared graph; Bazel and Vesta describe actions as content-addressed
functions of declared inputs~\cite{bazel,vesta}; \texttt{ccache} and \texttt{sccache} recognise a
compiler invocation and key on the source, the flags and the toolchain~\cite{ccache,sccache}; Nix
makes the derivation itself the unit of caching~\cite{nixos}. In each case the dependency set is
part of the interface. The system that does the caching and the system that knows the inputs are
the same system, or they agree on a description of them.

An arbitrary command has no such interface. A shell pipeline, a one-off compression step, a
scripted transformation --- none of them declare what they read. A cache for such commands must
either (a) discover the dependency set by observing an execution, or (b) assume determinism and
hope. Option (b) is cheap and is what several widely used tools effectively do, with the well-known
consequence that a program reading the clock or the environment can replay a stale answer. Option
(a) is what trace-based caching and process memoisation attempt; Rattle and related
record/replay systems demonstrated both the appeal and the difficulty~\cite{rr}.

This paper is about option (a) taken to its conservative extreme, and about what that costs. We
built \tracejit{} around a single rule: no cached byte is exposed unless every recorded dependency
has just been re-checked, and anything the tracer cannot re-check causes the command to run again.
We then measured the consequences, including the ones that make the system look bad.

\paragraph{The research question.}
\emph{If an automatic command cache refuses whenever it cannot prove a dependency is unchanged,
what does it cost in eligibility, latency and applicability, and which correctness properties does
that refusal actually buy?}

The question is deliberately two-sided. A paper that reported only the wins would be unable to say
whether the conservatism was worth anything; a paper that reported only the refusals would not show
a system at all. Our evaluation therefore keeps both, and the negative results are part of the
contribution rather than an appendix.

\paragraph{Contributions.}
\begin{enumerate}
  \item \textbf{A contract discipline for automatic reuse.} We separate the observable being
  reproduced (standard streams, exit status, declared outputs) from the evidence used to justify
  reuse (guards over recorded kernel observations), and we make the second a precondition of the
  first in the type system and in the runtime. Refusal, deoptimisation and re-execution are normal
  outcomes with their own reporting, not errors.
  \item \textbf{A measured account of the price of conservatism.} On a frozen census of 59
  workloads and three contracts, native eligibility is $4.1\%$ and audited-contract eligibility
  $24.5$--$38.8\%$; 28 of 162 attempted rows reused with zero incorrect observations. We report
  the refusals by cause, and we report the workloads that got slower.
  \item \textbf{A same-host competitive boundary.} Against \firebuild{} 0.8.2 and 0.8.6.2 we
  measure where each engine reuses, where only one does, and what each costs per session. The
  ordering flips with workload class, so we publish the boundary rather than a ranking.
  \item \textbf{Enforcement, qualified on a capable host.} We exercise the full lifecycle ---
  discovery, policy installation, \proven{} classification, reuse under enforcement, mutation,
  re-execution --- and report precisely which command class can complete it and which cannot.
  \item \textbf{The residual counterexample, kept.} A mutation between guard validation and output
  restoration produces a stale hit. We reproduce it, we explain why every economical userspace
  mitigation fails, and we state the operating assumption the prerelease requires.
\end{enumerate}

\paragraph{What we do not claim.}
We do not claim to have invented command caching, content-addressed storage, tracing, sandboxing,
or record/replay. We do not claim universal acceleration: several measured workloads are slower
with the cache than without it, and one compiler cache is more than an order of magnitude faster
than we are on its home ground. We do not claim formal correctness, and we do not describe our
enforcement label as a proof. Section~\ref{sec:negative} is where most of this is quantified.
